% !TeX root = ../DoAn.tex
\documentclass[../DoAn.tex]{subfiles}
\begin{document}

Tiếp nối các phân tích yêu cầu tại Chương 2 và nền tảng công nghệ ở Chương 3, Chương 4 trình bày thiết kế chi tiết, quá trình xây dựng, triển khai và đánh giá thực nghiệm hệ thống phần mềm tối ưu hóa quy trình điểm danh và quản lý học sinh. Nội dung chương tập trung vào việc mô tả thiết kế kiến trúc tổng thể của các phân hệ (\path{FE_WEB}, \path{frontend_student}, \path{frontend_teacher}, \path{backend-core}, \path{detect-face}, \path{ai-proctor}), thiết kế cơ sở dữ liệu PostgreSQL \cite{postgresql_docs2026}, xây dựng các kịch bản kiểm thử chức năng, và cấu hình triển khai thực tế bằng Docker Compose \cite{docker_compose_docs2026} kết hợp NGINX reverse proxy \cite{nginx_docs2026}.

\section{Thiết kế kiến trúc}
\label{section:4.1}
\subsection{Lựa chọn kiến trúc phần mềm}
\label{subsection:4.1.1}
Đối với hệ thống quản lý học tập tích hợp điểm danh Face ID và giám sát thi trực tuyến, đồ án lựa chọn kiến trúc phần mềm lai, trong đó phần nghiệp vụ trung tâm được tổ chức theo kiến trúc phân tầng ba lớp, còn các chức năng AI có tải xử lý riêng được tách thành các dịch vụ độc lập theo định hướng SOA/microservice nhẹ \cite{newman2021buildingmicroservices}. Cách lựa chọn này phù hợp với phạm vi đồ án vì hệ thống vừa cần giữ dữ liệu đào tạo nhất quán trong một backend chính, vừa cần tách các tác vụ xử lý ảnh và giám sát webcam ra khỏi luồng nghiệp vụ thông thường.

Về lý thuyết, kiến trúc phân tầng ba lớp chia hệ thống thành tầng giao diện, tầng xử lý nghiệp vụ và tầng dữ liệu. Khi áp dụng vào đồ án, các tầng này được ánh xạ thành những thành phần cụ thể như sau:
\begin{itemize}
    \item \textbf{Tầng giao diện}: gồm \path{FE_WEB} cho người quản trị, giảng viên và học viên trên trình duyệt; \path{frontend_student} cho học viên trên thiết bị di động; và \path{frontend_teacher} cho giảng viên trên thiết bị di động. Các client này không truy cập trực tiếp cơ sở dữ liệu mà gửi yêu cầu thông qua API do backend cung cấp.
    \item \textbf{Tầng xử lý nghiệp vụ}: là phân hệ \path{backend-core} được xây dựng bằng Spring Boot \cite{spring_boot_docs2026}. Bên trong backend, các gói \path{controller}, \path{service}, \path{repository}, \path{entity}, \path{dto} và \path{config} lần lượt đảm nhiệm vai trò tiếp nhận yêu cầu, xử lý nghiệp vụ, truy xuất dữ liệu, biểu diễn mô hình dữ liệu, trao đổi dữ liệu qua API và cấu hình hệ thống.
    \item \textbf{Tầng dữ liệu}: sử dụng PostgreSQL để lưu trữ các thực thể như người dùng, lớp học phần, bài thi, kết quả điểm danh, bài nộp và nhật ký vi phạm \cite{postgresql_docs2026}.
\end{itemize}

So với kiến trúc ba lớp thuần túy, hệ thống có bổ sung các dịch vụ chuyên biệt để xử lý nghiệp vụ AI:
\begin{itemize}
    \item \textbf{\path{detect-face}}: đảm nhiệm nhận diện khuôn mặt phục vụ điểm danh nhóm.
    \item \textbf{\path{ai-proctor}}: đảm nhiệm giám sát thi trực tuyến thông qua luồng webcam và WebSocket \cite{fastapi_docs2026,fette2011websocket}.
    \item \textbf{\path{backend-core}}: giữ vai trò điều phối, nhận yêu cầu từ client, kiểm tra điều kiện nghiệp vụ, gọi dịch vụ AI khi cần và lưu kết quả cuối cùng vào cơ sở dữ liệu.
\end{itemize}
Hai dịch vụ AI được triển khai độc lập với backend nghiệp vụ, giúp \path{backend-core} không phải mang theo các phụ thuộc nặng của xử lý ảnh, đồng thời cho phép cấu hình, thay thế hoặc mở rộng riêng khi tải xử lý AI tăng lên.

Như vậy, kiến trúc cụ thể của hệ thống có thể được tóm tắt theo ba nhóm thành phần chính:
\begin{itemize}
    \item \textbf{Nhóm client}: Web và ứng dụng di động.
    \item \textbf{Nhóm backend nghiệp vụ}: controller, service, repository, entity, DTO và cấu hình bảo mật.
    \item \textbf{Nhóm dịch vụ hạ tầng và AI}: PostgreSQL, Keycloak, \path{detect-face}, \path{ai-proctor}, NGINX và Docker Compose \cite{keycloak2026,nginx_docs2026,docker_compose_docs2026}.
\end{itemize}
Cách tổ chức này giữ được ưu điểm dễ phát triển, dễ kiểm soát dữ liệu của kiến trúc phân tầng, đồng thời bổ sung khả năng tách biệt và mở rộng cho các tác vụ AI đặc thù của đề tài.

\subsection{Thiết kế tổng quan}
\label{subsection:4.1.2}

Thiết kế tổng quan của hệ thống được tổ chức theo mô hình kiến trúc phân tầng, trong đó các thành phần được chia theo vai trò giao diện, xử lý nghiệp vụ, lưu trữ dữ liệu và dịch vụ AI chuyên biệt. Cách phân chia này giúp hệ thống dễ bảo trì, hạn chế phụ thuộc chéo và phù hợp với kiến trúc lai đã lựa chọn ở Mục \ref{subsection:4.1.1}. Các nguyên tắc thiết kế chính gồm:
\begin{itemize}
    \item \textbf{Phụ thuộc một chiều}: Các client chỉ gọi API; controller gọi service; service điều phối repository hoặc dịch vụ ngoài; tầng dữ liệu không phụ thuộc ngược lên tầng trên.
    \item \textbf{Không phụ thuộc vòng}: Các gói được tách theo trách nhiệm rõ ràng, tránh việc các module gọi qua lại lẫn nhau.
    \item \textbf{Tách riêng tác vụ AI}: Các chức năng nhận diện khuôn mặt và giám sát thi được triển khai độc lập để giảm tải cho backend nghiệp vụ.
\end{itemize}

Sơ đồ gói thiết kế tổng quan của hệ thống được trình bày trong Hình \ref{fig:BieuDoGoiThietKeTongQuan}.

\FloatBarrier
\begin{center}
    \centering
    \includegraphics[width=0.95\textwidth]{Hinhve/BieuDoGoiThietKeTongQuan.png}
    \captionof{figure}{Biểu đồ gói thiết kế tổng quan}
    \label{fig:BieuDoGoiThietKeTongQuan}
\end{center}
\FloatBarrier

Vai trò các tầng và gói chính được tóm tắt trong Bảng \ref{table:architecture_overview_packages}.

\begin{longtable}{|p{0.22\textwidth}|p{0.28\textwidth}|p{0.40\textwidth}|}
    \caption{Tóm tắt các tầng và gói chính trong thiết kế tổng quan}
    \label{table:architecture_overview_packages}\\
        \hline
        \textbf{Tầng/nhóm thành phần} & \textbf{Gói hoặc phân hệ chính} & \textbf{Vai trò} \\ \hline
    \endfirsthead
    \multicolumn{3}{r}{\textit{Tiếp theo trang trước}} \\
    \hline
    \textbf{Tầng/nhóm thành phần} & \textbf{Gói hoặc phân hệ chính} & \textbf{Vai trò} \\ \hline
    \endhead
    \hline
    \multicolumn{3}{r}{\textit{Xem tiếp trang sau}} \\
    \endfoot
    \hline
    \endlastfoot
        Giao diện người dùng & \path{FE_WEB}, \path{frontend_student}, \path{frontend_teacher} & Cung cấp giao diện Web và ứng dụng di động cho quản trị viên, giảng viên và học viên; gửi yêu cầu đến backend thông qua API. \\ \hline
        Nghiệp vụ và API & \path{controller}, \path{service}, \path{dto} & Tiếp nhận yêu cầu, kiểm tra dữ liệu đầu vào, xử lý nghiệp vụ điểm danh, thi trực tuyến, lớp học phần và trả phản hồi cho client. \\ \hline
        Truy xuất và mô hình dữ liệu & \path{repository}, \path{entity} & Ánh xạ dữ liệu nghiệp vụ sang các bảng PostgreSQL và thực hiện các thao tác truy xuất dữ liệu qua Spring Data JPA \cite{spring_boot_docs2026,postgresql_docs2026}. \\ \hline
        Hạ tầng dùng chung & \path{config}, \path{exception}, \path{util} & Cấu hình bảo mật, Keycloak SSO, CORS, WebSocket, xử lý ngoại lệ và các tiện ích dùng chung như JWT, GPS, file. \\ \hline
        Dịch vụ AI & \path{detect-face}, \path{ai-proctor} & Xử lý nhận diện khuôn mặt nhóm bằng RetinaFace/ArcFace và giám sát thi trực tuyến bằng webcam, WebSocket, MediaPipe FaceMesh và solvePnP \cite{Deng_2020_CVPR,Deng_2019_CVPR,Lugaresi2019MediaPipe,Lepetit2009EPnP}. \\ \hline
\end{longtable}

Nhìn chung, thiết kế tổng quan giữ backend Spring Boot làm trung tâm điều phối nghiệp vụ, PostgreSQL làm nơi lưu trữ dữ liệu chính, còn các tác vụ AI được tách thành dịch vụ độc lập. Cách tổ chức này giúp hệ thống duy trì ranh giới trách nhiệm rõ ràng giữa các tầng, đồng thời vẫn hỗ trợ mở rộng riêng các thành phần có tải xử lý cao.
\subsection{Thiết kế chi tiết gói}
\label{subsection:4.1.3}

Phần này trình bày thiết kế chi tiết gói cho hai phân hệ trọng tâm: \textbf{Quản lý thi trực tuyến \& Giám sát AI} và \textbf{Điểm danh thông minh FaceID nhóm \& Định vị GPS}. Mỗi phân hệ được mô tả bằng một sơ đồ gói và bảng tóm tắt vai trò các thành phần chính.

\subsubsection{Phân hệ Quản lý thi trực tuyến \& Giám sát phòng thi bằng trí tuệ nhân tạo (Assessment \& AI Proctoring)}
\label{subsubsection:4.1.3_assessment}

Phân hệ \textbf{Quản lý thi trực tuyến \& Giám sát AI (Assessment \& AI Proctoring)} đảm nhiệm các nghiệp vụ liên quan đến tổ chức bài kiểm tra, quản lý câu hỏi, ghi nhận bài làm và hỗ trợ chấm điểm. Bên cạnh đó, phân hệ còn tích hợp các thành phần giám sát bằng trí tuệ nhân tạo nhằm phát hiện những dấu hiệu bất thường trong quá trình làm bài, từ đó tăng tính minh bạch và độ tin cậy của hoạt động thi trực tuyến.

Chi tiết về cấu trúc các lớp (classes), giao diện (interfaces) và mối quan hệ giữa chúng trong phân hệ Quản lý thi \& Giám sát AI được thể hiện trong Hình \ref{fig:BieuDoThietKeChiTietGoi}.

\begin{figure}[!htbp]
    \centering
    \includegraphics[width=0.95\textwidth]{Hinhve/BieuDoThietKeChiTietGoi_QuanLyThi_GiamSatAI.png}
    \caption{Biểu đồ thiết kế chi tiết gói - Quản lý thi \& Giám sát AI}
    \label{fig:BieuDoThietKeChiTietGoi}
\end{figure}
\FloatBarrier

Sơ đồ trong Hình \ref{fig:BieuDoThietKeChiTietGoi} được tổ chức theo kiến trúc Spring Boot \cite{spring_boot_docs2026}, trong đó các gói chính được tóm tắt ở Bảng \ref{table:assessment_package_summary}.

\begin{longtable}{|p{0.22\textwidth}|p{0.68\textwidth}|}
    \caption{Tóm tắt các gói chính của phân hệ Quản lý thi \& Giám sát AI}
    \label{table:assessment_package_summary}\\
    \hline
    \textbf{Gói/thành phần} & \textbf{Vai trò chính} \\ \hline
    \endfirsthead
    \multicolumn{2}{r}{\textit{Tiếp theo trang trước}} \\
    \hline
    \textbf{Gói/thành phần} & \textbf{Vai trò chính} \\ \hline
    \endhead
    \hline
    \multicolumn{2}{r}{\textit{Xem tiếp trang sau}} \\
    \endfoot
    \hline
    \endlastfoot
    \path{controller} & Gồm \path{AssessmentController} và \path{AIProctorController}, tiếp nhận yêu cầu về đề thi, bài làm, cấu hình giám sát và nhật ký vi phạm. \\ \hline
    \path{dto} & Gồm \path{AssessmentDto}, \path{StudentSubmissionDto}, dùng để trao đổi dữ liệu giữa client và backend. \\ \hline
    \path{service}/\path{implement} & Gồm \path{AssessmentService} và \path{AssessmentServiceImplement}, xử lý nghiệp vụ soạn đề, lưu bài, chấm trắc nghiệm, kiểm soát thời gian và kết nối dịch vụ \path{ai-proctor}. \\ \hline
    \path{repository}/\path{entity} & Quản lý các thực thể \path{Assessment}, \path{AssessmentQuestion}, \path{StudentSubmission}, \path{StudentAnswer}, \path{Student} thông qua Spring Data JPA và PostgreSQL \cite{postgresql_docs2026}. \\ \hline
    \path{exception} & Chuẩn hóa lỗi nghiệp vụ như nộp bài quá giờ, truy cập không hợp lệ hoặc vi phạm do AI ghi nhận. \\ \hline
\end{longtable}

\subsubsection{Phân hệ Điểm danh thông minh bằng FaceID nhóm \& Định vị GPS (Smart Attendance)}
\label{subsubsection:4.1.3_attendance}

Bên cạnh phân hệ Quản lý thi, phân hệ \textbf{Điểm danh thông minh bằng FaceID nhóm \& Định vị GPS (Smart Attendance)} hỗ trợ giảng viên ghi nhận trạng thái tham gia học tập của học viên theo từng buổi học bằng cách kết hợp nhận diện khuôn mặt và kiểm tra vị trí GPS. Dữ liệu điểm danh được quản lý trong backend Spring Boot, lưu trữ bằng PostgreSQL \cite{spring_data_jpa_docs2026,postgresql_docs2026}; quá trình nhận diện khuôn mặt sử dụng RetinaFace và ArcFace \cite{Deng_2020_CVPR,Deng_2019_CVPR}; còn việc kiểm tra vị trí dựa trên công thức Haversine \cite{sinnott1984virtues}. Cách tách riêng các bước tiếp nhận dữ liệu, kiểm tra điều kiện và lưu kết quả giúp hệ thống dễ mở rộng, dễ bảo trì và thuận tiện cho việc đối soát khi có sai lệch.

Chi tiết về cấu trúc các lớp (classes), giao diện (interfaces) và mối quan hệ giữa chúng trong phân hệ Điểm danh thông minh được thể hiện trong Hình \ref{fig:BieuDoThietKeChiTietGoi_DiemDanh}.

\begin{figure}[!htbp]
    \centering
    \includegraphics[width=0.95\textwidth]{Hinhve/BieuDoThietKeChiTietGoi_DiemDanhFaceID_GPS.png}
    \caption{Biểu đồ thiết kế chi tiết gói - Điểm danh FaceID nhóm \& Định vị GPS}
    \label{fig:BieuDoThietKeChiTietGoi_DiemDanh}
\end{figure}
\FloatBarrier

Sơ đồ trong Hình \ref{fig:BieuDoThietKeChiTietGoi_DiemDanh} được tóm tắt theo các nhóm thành phần chính ở Bảng \ref{table:attendance_package_summary}.

\begin{longtable}{|p{0.22\textwidth}|p{0.68\textwidth}|}
    \caption{Tóm tắt các gói chính của phân hệ Điểm danh thông minh}
    \label{table:attendance_package_summary}\\
    \hline
    \textbf{Gói/thành phần} & \textbf{Vai trò chính} \\ \hline
    \endfirsthead
    \multicolumn{2}{r}{\textit{Tiếp theo trang trước}} \\
    \hline
    \textbf{Gói/thành phần} & \textbf{Vai trò chính} \\ \hline
    \endhead
    \hline
    \multicolumn{2}{r}{\textit{Xem tiếp trang sau}} \\
    \endfoot
    \hline
    \endlastfoot
    \path{controller} & Gồm \path{TeacherController} và \path{StudentController}, tiếp nhận yêu cầu tạo phiên, gửi GPS, điểm danh và xem kết quả. \\ \hline
    \path{dto} & Gồm \path{FormDto}, \path{StudentAttendancePreviewDto}, \path{ConfirmAttendanceChangesDto}, \path{LocationCheckRequest}, dùng để truyền dữ liệu điểm danh giữa client và backend. \\ \hline
    \path{service}/\path{implement} & Gồm \path{TeacherService}, \path{StudentService} và các lớp hiện thực, xử lý geofencing, chống giả lập GPS, kết nối \path{detect-face} và tổng hợp kết quả điểm danh. \\ \hline
    \path{repository}/\path{entity} & Quản lý \path{Teacher}, \path{Student}, \path{Register}, \path{Form}, \path{FormSubmission}, \path{AttendanceLog} bằng Spring Data JPA và PostgreSQL \cite{spring_data_jpa_docs2026,postgresql_docs2026}. \\ \hline
    \path{exception}/\path{spring-framework} & Chuẩn hóa phản hồi lỗi khi GPS ngoài phạm vi, FaceID không khớp hoặc phát hiện giả lập vị trí. \\ \hline
\end{longtable}

\section{Thiết kế chi tiết}
\subsection{Thiết kế giao diện}
Giao diện hệ thống được thiết kế theo hướng ứng dụng nghiệp vụ, ưu tiên thao tác nhanh, dễ đọc dữ liệu và hạn chế sai sót trong các quy trình như điểm danh, làm bài thi, chấm điểm và quản lý lớp học phần. Phần Web được hiện thực trong phân hệ \path{FE_WEB} bằng React, Vite và TailwindCSS \cite{react_docs2026,vite_docs2026,tailwindcss_docs2026}. Bố cục chính gồm \path{Navbar}, \path{Sidebar} và vùng nội dung trung tâm thay đổi theo vai trò người dùng.

Các chuẩn giao diện chính được tóm tắt trong Bảng \ref{table:ui_design_summary}. Hệ thống hướng tới màn hình máy tính từ $1366 \times 768$ đến $1920 \times 1080$ pixel cho thao tác nghiệp vụ thường xuyên; trên thiết bị nhỏ, \path{Sidebar} chuyển thành menu trượt, nội dung chính giảm khoảng đệm và các bảng/thẻ thông tin chuyển về bố cục một cột.

\begin{longtable}{|p{0.24\textwidth}|p{0.30\textwidth}|p{0.36\textwidth}|}
    \caption{Tóm tắt chuẩn thiết kế giao diện Web}
    \label{table:ui_design_summary}\\
        \hline
        \textbf{Nhóm thiết kế} & \textbf{Thành phần/giá trị sử dụng} & \textbf{Mục đích} \\ \hline
    \endfirsthead
    \multicolumn{3}{r}{\textit{Tiếp theo trang trước}} \\
    \hline
    \textbf{Nhóm thiết kế} & \textbf{Thành phần/giá trị sử dụng} & \textbf{Mục đích} \\ \hline
    \endhead
    \hline
    \multicolumn{3}{r}{\textit{Xem tiếp trang sau}} \\
    \endfoot
    \hline
    \endlastfoot
        Bố cục tổng thể & \path{Navbar}, \path{Sidebar}, vùng nội dung chính & Giữ khung thao tác thống nhất, giúp người dùng chuyển giữa các màn hình mà không phải học lại cách sử dụng. \\ \hline
        Hiển thị đáp ứng & Breakpoint và utility class của TailwindCSS & Tối ưu giao diện cho máy tính, máy tính bảng và điện thoại; menu trái có thể thu gọn trên màn hình nhỏ. \\ \hline
        Màu sắc và trạng thái & \path{primary}, \path{success}, \path{warning}, \path{error}, \path{info} & Phân biệt thao tác chính, trạng thái thành công, cảnh báo, lỗi và thông tin hệ thống. \\ \hline
        Kiểu chữ và thẻ nội dung & \path{Plus Jakarta Sans}, \path{Inter}, \path{Fira Code}, component \path{Card} & Tách rõ tiêu đề, nội dung, mã định danh và nhóm thông tin theo từng nhiệm vụ. \\ \hline
        Phản hồi thao tác & Loading, thông báo lỗi/thành công, vô hiệu hóa nút khi đang xử lý & Tránh gửi trùng yêu cầu và giúp người dùng biết trạng thái xử lý của hệ thống. \\ \hline
\end{longtable}

Giao diện được chia theo vai trò đăng nhập, không dùng một menu chung cho tất cả người dùng. Mỗi vai trò chỉ nhìn thấy các chức năng phù hợp với quyền hạn và trạng thái tài khoản của mình. Ví dụ, tài khoản giảng viên chưa được phê duyệt sẽ chuyển sang màn hình \path{PendingApproval}; học viên chưa hoàn thiện hồ sơ sẽ được yêu cầu bổ sung thông tin trước khi dùng các chức năng còn lại.

\begin{longtable}{|p{0.16\textwidth}|p{0.30\textwidth}|p{0.44\textwidth}|}
    \caption{Thiết kế giao diện theo vai trò người dùng}
    \label{table:ui_role_design}\\
        \hline
        \textbf{Vai trò} & \textbf{Màn hình/chức năng chính} & \textbf{Định hướng thiết kế} \\ \hline
    \endfirsthead
    \multicolumn{3}{r}{\textit{Tiếp theo trang trước}} \\
    \hline
    \textbf{Vai trò} & \textbf{Màn hình/chức năng chính} & \textbf{Định hướng thiết kế} \\ \hline
    \endhead
    \hline
    \multicolumn{3}{r}{\textit{Xem tiếp trang sau}} \\
    \endfoot
    \hline
    \endlastfoot
        Quản trị viên & \path{AdminDashboard}, \path{TeacherApproval}, \path{SemesterManagement} & Ưu tiên bảng dữ liệu, chỉ số tổng quan, nút phê duyệt/từ chối rõ ràng và trạng thái xử lý minh bạch. \\ \hline
        Giảng viên & \path{ClassManagement}, \path{ClassDetail}, \path{AssessmentManagement}, \path{GradeAssessment}, \path{PhotoAttendance} & Hỗ trợ quản lý lớp học phần, tài liệu, bài kiểm tra, chấm điểm và điểm danh; riêng điểm danh ảnh có bước xem trước trước khi lưu. \\ \hline
        Học viên & \path{MyCourses}, \path{TakeAssessment}, \path{GradesAndAttendance}, \path{Timetable}, \path{FaceUpload} & Giảm số bước thao tác, tập trung vào xem lớp học, tài liệu, lịch học, làm bài thi và tra cứu điểm/chuyên cần. \\ \hline
\end{longtable}

Về tương tác, các thao tác ảnh hưởng đến dữ liệu quan trọng đều có bước kiểm tra trước khi ghi nhận. Với điểm danh bằng ảnh lớp học, hệ thống tạo danh sách xem trước để giảng viên đối soát rồi mới xác nhận lưu. Với bài thi trực tuyến, hệ thống kiểm soát thời gian mở/đóng và chỉ cho phép nộp khi thỏa mãn điều kiện nghiệp vụ. Cách thiết kế này giúp giao diện vừa gọn, vừa phù hợp với đặc thù của hệ thống quản lý học tập tích hợp điểm danh và giám sát AI.
\subsection{Thiết kế lớp}
Thiết kế lớp đặc tả cấu trúc của các thuộc tính và phương thức trong những lớp đối tượng chủ đạo, đồng thời mô tả luồng thông điệp động qua biểu đồ trình tự để giải quyết hai nghiệp vụ cốt lõi: thi trực tuyến giám sát AI và điểm danh FaceID nhóm kết hợp định vị GPS.

\subsubsection{Đặc tả cấu trúc các lớp chính}
Hệ thống sử dụng các lớp cốt lõi được tổ chức theo kiến trúc phân tầng của Spring Boot. Nhóm \texttt{entity} biểu diễn dữ liệu lưu trong PostgreSQL, nhóm \texttt{repository} truy xuất dữ liệu, nhóm \texttt{service}/\texttt{implement} xử lý nghiệp vụ, còn nhóm \texttt{dto} chuẩn hóa dữ liệu trao đổi giữa client và backend. Trong phạm vi đồ án, phần thiết kế lớp tập trung vào hai nghiệp vụ chính: thi trực tuyến có giám sát AI và điểm danh thông minh bằng FaceID/GPS.

\begin{itemize}
    \item \textbf{Lớp thực thể \texttt{Assessment}}: Đại diện cho bài thi, lưu thông tin như tên bài thi, thời gian mở/đóng, danh sách câu hỏi và các cấu hình kiểm soát như \texttt{isLocationRequired}, \texttt{isCameraRequired}, \texttt{allowedRadiusMeters}, \texttt{teacherLatitude}, \texttt{teacherLongitude}. Các thuộc tính này cho phép từng bài thi tự xác định yêu cầu GPS, camera và bán kính làm bài hợp lệ.
    
    \item \textbf{Nhóm lớp bài làm}: Gồm \texttt{StudentSubmission} và \texttt{StudentAnswer}. \texttt{StudentSubmission} lưu một lần làm bài của học viên, bao gồm trạng thái làm bài, thời điểm bắt đầu, thời điểm nộp và tổng điểm; \texttt{StudentAnswer} lưu câu trả lời cụ thể cho từng câu hỏi. Nhóm lớp này giúp tách dữ liệu đề thi khỏi dữ liệu bài làm, thuận tiện cho lưu nháp, nộp bài và chấm điểm tự động.
    
    \item \textbf{Lớp nghiệp vụ \texttt{AssessmentServiceImplement}}: Điều phối nghiệp vụ thi trực tuyến giữa controller, repository và dịch vụ giám sát AI. Các phương thức chính gồm:
    \begin{itemize}
        \item \texttt{startAssessment(assessmentId, studentId, location)}: Khởi tạo bài thi, gọi \texttt{validateAssessmentLocation} để kiểm tra GPS và lưu trạng thái \texttt{IN\_PROGRESS}.
        \item \texttt{submitAssessment(submissionId, studentId, location)}: Xác thực vị trí lần cuối, tự động chấm điểm trắc nghiệm và câu trả lời ngắn qua \texttt{autoGradeMultipleChoice} và \texttt{autoGradeShortAnswer}, cập nhật trạng thái bài thi thành \texttt{GRADED} hoặc \texttt{SUBMITTED}.
        \item \texttt{validateAssessmentLocation(assessment, location)}: Áp dụng công thức Haversine \cite{sinnott1984virtues} để tính khoảng cách địa lý thực tế và quét cờ giả lập vị trí (\texttt{mockLocationDetected}).
    \end{itemize}
    
    \item \textbf{Nhóm lớp điểm danh}: Gồm \texttt{Form}, \texttt{FormSubmission}, \texttt{Register} và \texttt{AttendanceLog}. \texttt{Form} biểu diễn phiên điểm danh, chứa mã OTP, thời gian mở và cấu hình vị trí; \texttt{FormSubmission} lưu từng lượt check-in; \texttt{Register} lưu hồ sơ FaceID gốc; còn \texttt{AttendanceLog} lưu lịch sử chuyên cần sau khi kết quả được xác nhận.
    
    \item \textbf{Lớp nghiệp vụ \texttt{TeacherServiceImplement}}: Thực thi nghiệp vụ điểm danh phía giảng viên, từ tạo phiên điểm danh, nhận kết quả nhận diện từ dịch vụ \texttt{detect-face}, tạo danh sách xem trước đến lưu kết quả chính thức. Các phương thức chính gồm:
    \begin{itemize}
        \item \texttt{createForm(courseId, formDto, sub)}: Tạo phiên điểm danh kèm mã OTP và cấu hình Geofencing.
        \item \texttt{previewAttendanceFace(courseId, lectureNumber, recognizedStudentIds, sub)}: Nhận danh sách ID sinh viên được AI phát hiện từ ảnh tập thể để chuẩn bị bảng xem trước đối sánh.
        \item \texttt{confirmAttendanceChanges(dto, sub)}: Giảng viên duyệt chỉnh sửa và lưu chính thức kết quả điểm danh vào \texttt{AttendanceLog}, đồng thời chạy stored procedure tự động cập nhật số liệu chuyên cần.
    \end{itemize}
    
    \item \textbf{Lớp nghiệp vụ \texttt{StudentServiceImplement}}: Xử lý các thao tác phía học viên như đăng ký hồ sơ khuôn mặt, nhập OTP, gửi tọa độ GPS và check-in cá nhân. Lớp này phối hợp với các repository liên quan để xác minh điều kiện điểm danh trước khi tạo bản ghi hợp lệ.
\end{itemize}

\subsubsection{Biểu đồ trình tự mô tả luồng truyền thông điệp}

\paragraph{1. UC-10: Làm bài thi trực tuyến và Giám sát AI}
Biểu đồ trình tự mô tả quy trình học viên làm bài thi trực tuyến dưới sự giám sát liên tục của dịch vụ AI thời gian thực và ghi nhận vi phạm về máy chủ Backend, được minh họa ở Hình \ref{fig:UC10_sequence_diagram}.

\begin{figure}[H]
    \centering
    \includegraphics[width=0.9\textwidth]{Hinhve/UC_10_Thi_truc_tuyen_Giam_sat_AI.png}
    \caption{Biểu đồ trình tự UC-10: Làm bài thi trực tuyến và Giám sát AI}
    \label{fig:UC10_sequence_diagram}
\end{figure}
\FloatBarrier

Quy trình tương tác diễn ra như sau:
\begin{enumerate}
    \item \textbf{Khởi tạo thi}: Client (\texttt{FE\_WEB / App}) lấy đề qua \texttt{getAssessmentsForCourse} và gọi API \texttt{startAssessment}. Lớp \texttt{AssessmentService} gọi \texttt{validateAssessmentLocation} xác thực GPS và lưu bài nộp \texttt{StudentSubmission} trạng thái \texttt{IN\_PROGRESS}.
    \item \textbf{Giám sát AI trực tuyến}: Client truyền luồng webcam qua WebSocket tới \texttt{ai-proctor}. Dịch vụ AI dùng MediaPipe FaceMesh \cite{Lugaresi2019MediaPipe} và EPnP \cite{Lepetit2009EPnP} ước lượng hướng nhìn. Nếu có bất thường trên 3 giây, hệ thống lưu video vi phạm 10 giây và gửi cảnh báo về backend.
    \item \textbf{Nộp bài}: Client gửi yêu cầu \texttt{submitAssessment} kèm GPS. Hệ thống tự động chấm điểm trắc nghiệm, điền khuyết và lưu kết quả với trạng thái \texttt{GRADED} hoặc \texttt{SUBMITTED}.
\end{enumerate}

\paragraph{2. UC-06: Điểm danh nhận diện FaceID nhóm và Định vị GPS}
Biểu đồ trình tự mô tả quy trình giảng viên điểm danh nhóm bằng AI FaceID kết hợp với học viên check-in cá nhân bằng định vị GPS phòng học và mã OTP bổ sung, được minh họa ở Hình \ref{fig:UC06_sequence_diagram}.

\begin{figure}[H]
    \centering
    \includegraphics[width=0.9\textwidth]{Hinhve/UC_06_Diem_danh_FaceID_nhom_GPS.png}
    \caption{Biểu đồ trình tự UC-06: Điểm danh nhận diện FaceID nhóm và Định vị GPS}
    \label{fig:UC06_sequence_diagram}
\end{figure}
\FloatBarrier

Quy trình tương tác diễn ra như sau:
\begin{enumerate}
    \item \textbf{Điểm danh nhóm FaceID}: Giảng viên gửi ảnh lớp học tới AI Service (\texttt{detect-face}). Dịch vụ sử dụng RetinaFace \cite{Deng_2020_CVPR} và ArcFace \cite{Deng_2019_CVPR} để phát hiện và đối sánh khuôn mặt để xác định học viên. Kết quả gửi về \texttt{TeacherController} dưới dạng danh sách xem trước. Giảng viên đối soát và lưu qua \texttt{confirmAttendanceChanges}, hệ thống tự động cập nhật số buổi học của học viên.
    \item \textbf{Check-in cá nhân}: Giảng viên cấp mã OTP. Học viên nhập OTP, chụp ảnh selfie (xác thực qua ArcFace) và gửi tọa độ GPS. Lớp nghiệp vụ \texttt{StudentService} xác thực vị trí bằng Geofencing, kiểm tra chống giả lập GPS và lưu kết quả vào bảng \texttt{FormSubmission}.
\end{enumerate}
\subsection{Thiết kế cơ sở dữ liệu}
\label{subsection:4.2.3}

Hệ thống sử dụng hệ quản trị cơ sở dữ liệu PostgreSQL \cite{postgresql_docs2026} để lưu trữ thông tin nghiệp vụ và hỗ trợ tính toán khoảng cách GPS thông qua kiểu dữ liệu số thực chính xác kép (\path{double precision}). Các bảng dữ liệu được thiết kế chuẩn hóa, ràng buộc khóa ngoại chặt chẽ nhằm bảo đảm tính toàn vẹn và hạn chế dư thừa dữ liệu.

\subsubsection{Biểu đồ quan hệ thực thể (ERD)}
Mối quan hệ giữa các thực thể cốt lõi trong hệ thống được thể hiện chi tiết qua Biểu đồ quan hệ thực thể (E-R Diagram) trong Hình \ref{fig:BieuDoQuanHeThucThe}.

\FloatBarrier
\begin{center}
    \centering
    \includegraphics[width=0.95\textwidth]{Hinhve/BieuDoQuanHeThucThe.png}
    \captionof{figure}{Biểu đồ quan hệ thực thể (E-R) của hệ thống}
    \label{fig:BieuDoQuanHeThucThe}
\end{center}
\FloatBarrier

Sơ đồ cơ sở dữ liệu gồm ba nhóm thực thể chính:
\begin{itemize}
    \item \textbf{Nhóm quản lý người dùng và đào tạo}: Gồm các bảng \path{student}, \path{teacher}, \path{semester}, \path{course} và \path{register}, liên kết định danh qua trường \path{keycloak_id}.
    \item \textbf{Nhóm chuyên cần và điểm danh}: Gồm các bảng \path{attendance_log}, \path{form}, \path{form_submission}, \path{question} và \path{answer}. Bảng \path{form} cấu hình định vị lớp học, \path{form_submission} lưu chi tiết check-in (tọa độ, cờ giả lập và FaceID), sau đó tổng hợp vào \path{attendance_log}.
    \item \textbf{Nhóm quản lý thi và giám sát AI}: Gồm các bảng \path{assessment}, \path{assessment_question}, \path{student_submission} và \path{student_answer}. Bảng \path{assessment} lưu cấu hình thi (GPS/camera), \path{student_submission} lưu thông tin làm bài (tọa độ, cờ giả lập) và \path{student_answer} lưu chi tiết câu trả lời.
\end{itemize}

\subsubsection{Chi tiết các bảng dữ liệu chính}
Dưới đây là đặc tả chi tiết cấu trúc các cột, kiểu dữ liệu và ràng buộc của các bảng dữ liệu cốt lõi trong hệ thống.

\paragraph{1. Bảng student (Học viên)}
Bảng này lưu trữ thông tin cá nhân của học viên và đường dẫn đến ảnh khuôn mẫu gốc phục vụ đối sánh FaceID.
\begin{longtable}{|p{0.25\textwidth}|p{0.22\textwidth}|p{0.46\textwidth}|}
    \caption{Đặc tả cấu trúc bảng student}
    \label{table:db_student}\\
    \hline
    \textbf{Tên trường} & \textbf{Kiểu dữ liệu} & \textbf{Ý nghĩa/Ràng buộc} \\ \hline
    \endfirsthead
    \multicolumn{3}{r}{\textit{Tiếp theo trang trước}} \\
    \hline
    \textbf{Tên trường} & \textbf{Kiểu dữ liệu} & \textbf{Ý nghĩa/Ràng buộc} \\ \hline
    \endhead
    \hline
    \multicolumn{3}{r}{\textit{Xem tiếp trang sau}} \\
    \endfoot
    \hline
    \endlastfoot
    \path{id} & \path{bigint} & Khóa chính, tự động tăng. \\ \hline
    \path{student_code} & \path{varchar} & Mã số sinh viên (Duy nhất, không null). \\ \hline
    \path{name} & \path{varchar} & Họ và tên học viên. \\ \hline
    \path{email} & \path{varchar} & Địa chỉ email của học viên. \\ \hline
    \path{keycloak_id} & \path{varchar} & ID định danh của tài khoản trên hệ thống Keycloak SSO. \\ \hline
    \path{image_path} & \path{varchar} & Đường dẫn tệp ảnh FaceID mẫu của học viên lưu trên server. \\ \hline
    \path{is_active} & \path{boolean} & Trạng thái kích hoạt tài khoản. \\ \hline
    \path{created_at} & \path{timestamp} & Thời điểm tạo bản ghi. \\ \hline
    \path{updated_at} & \path{timestamp} & Thời điểm cập nhật bản ghi gần nhất. \\ \hline
\end{longtable}

\paragraph{2. Bảng teacher (Giảng viên)}
Bảng này lưu trữ thông tin giảng viên và trạng thái phê duyệt tài khoản của hệ thống quản trị.
\begin{longtable}{|p{0.25\textwidth}|p{0.22\textwidth}|p{0.46\textwidth}|}
    \caption{Đặc tả cấu trúc bảng teacher}
    \label{table:db_teacher}\\
    \hline
    \textbf{Tên trường} & \textbf{Kiểu dữ liệu} & \textbf{Ý nghĩa/Ràng buộc} \\ \hline
    \endfirsthead
    \multicolumn{3}{r}{\textit{Tiếp theo trang trước}} \\
    \hline
    \textbf{Tên trường} & \textbf{Kiểu dữ liệu} & \textbf{Ý nghĩa/Ràng buộc} \\ \hline
    \endhead
    \hline
    \multicolumn{3}{r}{\textit{Xem tiếp trang sau}} \\
    \endfoot
    \hline
    \endlastfoot
    \path{id} & \path{bigint} & Khóa chính, tự động tăng. \\ \hline
    \path{teacher_code} & \path{varchar} & Mã giảng viên (Duy nhất, không null). \\ \hline
    \path{name} & \path{varchar} & Họ và tên giảng viên. \\ \hline
    \path{email} & \path{varchar} & Địa chỉ email của giảng viên. \\ \hline
    \path{keycloak_id} & \path{varchar} & ID định danh của tài khoản trên Keycloak SSO. \\ \hline
    \path{account_status} & \path{varchar} & Trạng thái phê duyệt (\path{PENDING}, \path{APPROVED}, \path{REJECTED}). \\ \hline
    \path{rejection_reason} & \path{varchar} & Lý do từ chối tài khoản (nếu trạng thái là \path{REJECTED}). \\ \hline
    \path{is_active} & \path{boolean} & Trạng thái hoạt động. \\ \hline
    \path{created_at} & \path{timestamp} & Thời điểm tạo bản ghi. \\ \hline
    \path{updated_at} & \path{timestamp} & Thời điểm cập nhật gần nhất. \\ \hline
\end{longtable}

\paragraph{3. Bảng course (Lớp học phần)}
Lớp học phần liên kết giữa môn học, giảng viên phụ trách và học kỳ tương ứng.
\begin{longtable}{|p{0.25\textwidth}|p{0.22\textwidth}|p{0.46\textwidth}|}
    \caption{Đặc tả cấu trúc bảng course}
    \label{table:db_course}\\
    \hline
    \textbf{Tên trường} & \textbf{Kiểu dữ liệu} & \textbf{Ý nghĩa/Ràng buộc} \\ \hline
    \endfirsthead
    \multicolumn{3}{r}{\textit{Tiếp theo trang trước}} \\
    \hline
    \textbf{Tên trường} & \textbf{Kiểu dữ liệu} & \textbf{Ý nghĩa/Ràng buộc} \\ \hline
    \endhead
    \hline
    \multicolumn{3}{r}{\textit{Xem tiếp trang sau}} \\
    \endfoot
    \hline
    \endlastfoot
    \path{id} & \path{bigint} & Khóa chính, tự động tăng. \\ \hline
    \path{course_code} & \path{varchar} & Mã lớp học phần (Ví dụ: INT3306 1, không null). \\ \hline
    \path{subject} & \path{varchar} & Tên môn học học phần. \\ \hline
    \path{description} & \path{varchar} & Mô tả chi tiết về lớp học phần. \\ \hline
    \path{teacher_id} & \path{bigint} & Khóa ngoại liên kết bảng \path{teacher(id)}. \\ \hline
    \path{semester_id} & \path{bigint} & Khóa ngoại liên kết bảng \path{semester(id)}. \\ \hline
    \path{is_active} & \path{boolean} & Trạng thái hoạt động của lớp học phần. \\ \hline
    \path{created_at} & \path{timestamp} & Thời điểm tạo bản ghi. \\ \hline
    \path{updated_at} & \path{timestamp} & Thời điểm cập nhật gần nhất. \\ \hline
\end{longtable}

\paragraph{4. Bảng register (Đăng ký lớp học phần)}
Bảng liên kết nhiều-nhiều giữa học viên và lớp học phần, đồng thời lưu trữ tổng hợp số liệu chuyên cần để phục vụ đối soát nhanh.
\begin{longtable}{|p{0.25\textwidth}|p{0.22\textwidth}|p{0.46\textwidth}|}
    \caption{Đặc tả cấu trúc bảng register}
    \label{table:db_register}\\
    \hline
    \textbf{Tên trường} & \textbf{Kiểu dữ liệu} & \textbf{Ý nghĩa/Ràng buộc} \\ \hline
    \endfirsthead
    \multicolumn{3}{r}{\textit{Tiếp theo trang trước}} \\
    \hline
    \textbf{Tên trường} & \textbf{Kiểu dữ liệu} & \textbf{Ý nghĩa/Ràng buộc} \\ \hline
    \endhead
    \hline
    \multicolumn{3}{r}{\textit{Xem tiếp trang sau}} \\
    \endfoot
    \hline
    \endlastfoot
    \path{student_id} & \path{bigint} & Khóa chính \& Khóa ngoại liên kết bảng \path{student(id)}. \\ \hline
    \path{course_id} & \path{bigint} & Khóa chính \& Khóa ngoại liên kết bảng \path{course(id)}. \\ \hline
    \path{register_time} & \path{timestamp} & Thời điểm học viên đăng ký vào lớp học phần. \\ \hline
    \path{note} & \path{varchar} & Ghi chú thêm. \\ \hline
    \path{number_of_attendance} & \path{integer} & Tổng số buổi học viên đi học thực tế (Được cập nhật tự động). \\ \hline
    \path{number_of_absence} & \path{integer} & Tổng số buổi học viên vắng mặt (Được cập nhật tự động). \\ \hline
    \path{can_upload_documents} & \path{boolean} & Quyền tải tài liệu lên thư mục học phần. \\ \hline
    \path{can_download_documents} & \path{boolean} & Quyền tải tài liệu từ thư mục học phần. \\ \hline
\end{longtable}

\paragraph{5. Bảng attendance\_log (Nhật ký chuyên cần)}
Lưu kết quả điểm danh chính thức của từng học viên cho từng buổi học (sau khi giảng viên đã đối soát và phê duyệt).
\begin{longtable}{|p{0.25\textwidth}|p{0.22\textwidth}|p{0.46\textwidth}|}
    \caption{Đặc tả cấu trúc bảng attendance\_log}
    \label{table:db_attendance_log}\\
    \hline
    \textbf{Tên trường} & \textbf{Kiểu dữ liệu} & \textbf{Ý nghĩa/Ràng buộc} \\ \hline
    \endfirsthead
    \multicolumn{3}{r}{\textit{Tiếp theo trang trước}} \\
    \hline
    \textbf{Tên trường} & \textbf{Kiểu dữ liệu} & \textbf{Ý nghĩa/Ràng buộc} \\ \hline
    \endhead
    \hline
    \multicolumn{3}{r}{\textit{Xem tiếp trang sau}} \\
    \endfoot
    \hline
    \endlastfoot
    \path{id} & \path{bigint} & Khóa chính, tự động tăng. \\ \hline
    \path{student_id} & \path{bigint} & Khóa ngoại liên kết bảng \path{student(id)}. \\ \hline
    \path{course_id} & \path{bigint} & Khóa ngoại liên kết bảng \path{course(id)}. \\ \hline
    \path{lecture_number} & \path{integer} & Số thứ tự buổi học của học phần (Ví dụ: Buổi 1, 2, ...). \\ \hline
    \path{attendance_time} & \path{timestamp} & Thời gian chốt chuyên cần của buổi học. \\ \hline
    \path{is_attendance} & \path{boolean} & Kết quả điểm danh (\path{true}: Có mặt, \path{false}: Vắng mặt). \\ \hline
\end{longtable}

\paragraph{6. Bảng form (Phiên điểm danh)}
Cấu hình chi tiết một phiên điểm danh do giảng viên tạo, chứa mã OTP và các quy định về vị trí GPS và sinh trắc học.
\begin{longtable}{|p{0.25\textwidth}|p{0.22\textwidth}|p{0.46\textwidth}|}
    \caption{Đặc tả cấu trúc bảng form}
    \label{table:db_form}\\
    \hline
    \textbf{Tên trường} & \textbf{Kiểu dữ liệu} & \textbf{Ý nghĩa/Ràng buộc} \\ \hline
    \endfirsthead
    \multicolumn{3}{r}{\textit{Tiếp theo trang trước}} \\
    \hline
    \textbf{Tên trường} & \textbf{Kiểu dữ liệu} & \textbf{Ý nghĩa/Ràng buộc} \\ \hline
    \endhead
    \hline
    \multicolumn{3}{r}{\textit{Xem tiếp trang sau}} \\
    \endfoot
    \hline
    \endlastfoot
    \path{id} & \path{bigint} & Khóa chính, tự động tăng. \\ \hline
    \path{course_id} & \path{bigint} & Khóa ngoại liên kết bảng \path{course(id)}. \\ \hline
    \path{code} & \path{varchar} & Mã OTP điểm danh (giảng viên công bố). \\ \hline
    \path{lecture_number} & \path{integer} & Số thứ tự buổi học của phiên điểm danh này. \\ \hline
    \path{latitude} & \path{double} & Vĩ độ của vị trí lớp học do giảng viên thiết lập khi tạo form. \\ \hline
    \path{longitude} & \path{double} & Kinh độ của vị trí lớp học do giảng viên thiết lập khi tạo form. \\ \hline
    \path{is_location_required} & \path{boolean} & Cờ yêu cầu kiểm tra GPS Geofencing khi check-in. \\ \hline
    \path{allowed_radius_meters} & \path{integer} & Bán kính GPS tối đa cho phép để check-in hợp lệ (mét). \\ \hline
    \path{is_face_verification_required} & \path{boolean} & Cờ yêu cầu xác thực FaceID cá nhân khi check-in. \\ \hline
    \path{expired_at} & \path{timestamp} & Thời điểm phiên điểm danh hết hạn. \\ \hline
    \path{created_at} & \path{timestamp} & Thời điểm tạo phiên điểm danh. \\ \hline
\end{longtable}

\paragraph{7. Bảng form\_submission (Lượt check-in điểm danh)}
Lưu trữ thông tin chi tiết từng lượt gửi check-in điểm danh từ thiết bị di động của học viên để đối soát điều kiện GPS và FaceID.
\begin{longtable}{|p{0.25\textwidth}|p{0.22\textwidth}|p{0.46\textwidth}|}
    \caption{Đặc tả cấu trúc bảng form\_submission}
    \label{table:db_form_submission}\\
    \hline
    \textbf{Tên trường} & \textbf{Kiểu dữ liệu} & \textbf{Ý nghĩa/Ràng buộc} \\ \hline
    \endfirsthead
    \multicolumn{3}{r}{\textit{Tiếp theo trang trước}} \\
    \hline
    \textbf{Tên trường} & \textbf{Kiểu dữ liệu} & \textbf{Ý nghĩa/Ràng buộc} \\ \hline
    \endhead
    \hline
    \multicolumn{3}{r}{\textit{Xem tiếp trang sau}} \\
    \endfoot
    \hline
    \endlastfoot
    \path{id} & \path{bigint} & Khóa chính, tự động tăng. \\ \hline
    \path{student_id} & \path{bigint} & Khóa ngoại liên kết bảng \path{student(id)}. \\ \hline
    \path{form_id} & \path{bigint} & Khóa ngoại liên kết bảng \path{form(id)}. \\ \hline
    \path{is_correct} & \path{boolean} & Xác nhận học viên nhập đúng mã OTP điểm danh. \\ \hline
    \path{student_latitude} & \path{double} & Vĩ độ GPS thực tế của thiết bị học viên khi check-in. \\ \hline
    \path{student_longitude} & \path{double} & Kinh độ GPS thực tế của thiết bị học viên khi check-in. \\ \hline
    \path{calculated_distance} & \path{double} & Khoảng cách từ học viên tới lớp học tính bằng Haversine \cite{sinnott1984virtues}. \\ \hline
    \path{is_valid_location} & \path{boolean} & Xác nhận tọa độ hợp lệ (khoảng cách nhỏ hơn bán kính cấu hình). \\ \hline
    \path{mock_location_detected} & \path{boolean} & Cờ ghi nhận việc phát hiện ứng dụng giả lập GPS trên thiết bị. \\ \hline
    \path{is_face_verified} & \path{boolean} & Xác nhận khuôn mặt selfie trùng khớp ảnh mẫu sinh trắc học gốc. \\ \hline
    \path{submitted_at} & \path{timestamp} & Thời điểm thực hiện check-in. \\ \hline
\end{longtable}

\paragraph{8. Bảng assessment (Đề thi / Bài tập)}
Cấu hình đề thi trắc nghiệm/tự luận, thời gian thi, và quy chế giám sát thi trực tuyến bằng camera/AI.
\begin{longtable}{|p{0.25\textwidth}|p{0.22\textwidth}|p{0.46\textwidth}|}
    \caption{Đặc tả cấu trúc bảng assessment}
    \label{table:db_assessment}\\
    \hline
    \textbf{Tên trường} & \textbf{Kiểu dữ liệu} & \textbf{Ý nghĩa/Ràng buộc} \\ \hline
    \endfirsthead
    \multicolumn{3}{r}{\textit{Tiếp theo trang trước}} \\
    \hline
    \textbf{Tên trường} & \textbf{Kiểu dữ liệu} & \textbf{Ý nghĩa/Ràng buộc} \\ \hline
    \endhead
    \hline
    \multicolumn{3}{r}{\textit{Xem tiếp trang sau}} \\
    \endfoot
    \hline
    \endlastfoot
    \path{id} & \path{bigint} & Khóa chính, tự động tăng. \\ \hline
    \path{course_id} & \path{bigint} & Khóa ngoại liên kết bảng \path{course(id)}. \\ \hline
    \path{title} & \path{varchar} & Tiêu đề bài thi/bài tập (Không null). \\ \hline
    \path{description} & \path{varchar} & Mô tả và hướng dẫn làm bài. \\ \hline
    \path{type} & \path{varchar} & Loại bài (\path{EXAM}: Thi/Kiểm tra, \path{ASSIGNMENT}: Bài tập). \\ \hline
    \path{max_score} & \path{double} & Điểm số tối đa của bài thi. \\ \hline
    \path{duration_minutes} & \path{integer} & Thời lượng làm bài (tính bằng phút). \\ \hline
    \path{deadline} & \path{timestamp} & Hạn chót cho phép làm bài và nộp bài. \\ \hline
    \path{score_release_mode} & \path{varchar} & Chế độ công bố điểm (\path{AUTO}: Tự động, \path{MANUAL}: Thủ công). \\ \hline
    \path{is_published} & \path{boolean} & Trạng thái đã xuất bản đề thi tới học viên. \\ \hline
    \path{is_location_required} & \path{boolean} & Cờ bắt buộc kiểm tra GPS phòng thi khi làm bài. \\ \hline
    \path{is_camera_required} & \path{boolean} & Cờ bắt buộc mở camera để AI giám sát trong suốt quá trình. \\ \hline
    \path{allowed_radius_meters} & \path{integer} & Bán kính giới hạn địa lý của phòng thi (mét). \\ \hline
    \path{teacher_latitude} & \path{double} & Vĩ độ tâm của phòng thi (Geofencing). \\ \hline
    \path{teacher_longitude} & \path{double} & Kinh độ tâm của phòng thi (Geofencing). \\ \hline
    \path{weight} & \path{double} & Trọng số của điểm thi trong điểm học phần (từ 0.0 đến 1.0). \\ \hline
    \path{created_at} & \path{timestamp} & Thời điểm tạo đề thi. \\ \hline
    \path{updated_at} & \path{timestamp} & Thời điểm chỉnh sửa gần nhất. \\ \hline
\end{longtable}

\paragraph{9. Bảng student\_submission (Bài nộp thi học viên)}
Lưu trữ trạng thái làm bài, điểm số cuối cùng, tọa độ làm bài thi, và ghi nhận gian lận của học viên khi tham gia thi trực tuyến.
\begin{longtable}{|p{0.25\textwidth}|p{0.22\textwidth}|p{0.46\textwidth}|}
    \caption{Đặc tả cấu trúc bảng student\_submission}
    \label{table:db_student_submission}\\
    \hline
    \textbf{Tên trường} & \textbf{Kiểu dữ liệu} & \textbf{Ý nghĩa/Ràng buộc} \\ \hline
    \endfirsthead
    \multicolumn{3}{r}{\textit{Tiếp theo trang trước}} \\
    \hline
    \textbf{Tên trường} & \textbf{Kiểu dữ liệu} & \textbf{Ý nghĩa/Ràng buộc} \\ \hline
    \endhead
    \hline
    \multicolumn{3}{r}{\textit{Xem tiếp trang sau}} \\
    \endfoot
    \hline
    \endlastfoot
    \path{id} & \path{bigint} & Khóa chính, tự động tăng. \\ \hline
    \path{assessment_id} & \path{bigint} & Khóa ngoại liên kết bảng \path{assessment(id)}. \\ \hline
    \path{student_id} & \path{varchar} & ID tài khoản học viên (Liên kết Keycloak ID, không null). \\ \hline
    \path{status} & \path{varchar} & Trạng thái bài làm (\path{IN_PROGRESS}, \path{SUBMITTED}, \path{GRADED}). \\ \hline
    \path{final_score} & \path{double} & Điểm số cuối cùng đạt được sau khi chấm bài. \\ \hline
    \path{teacher_feedback} & \path{varchar} & Nhận xét, đánh giá của giảng viên chấm thi. \\ \hline
    \path{student_latitude} & \path{double} & Vĩ độ GPS thực tế của học viên khi thực hiện làm bài. \\ \hline
    \path{student_longitude} & \path{double} & Kinh độ GPS thực tế của học viên khi thực hiện làm bài. \\ \hline
    \path{calculated_distance} & \path{double} & Khoảng cách thực tế đến tọa độ phòng thi cấu hình. \\ \hline
    \path{is_valid_location} & \path{boolean} & Xác nhận học viên thi đúng vị trí phòng thi hợp lệ. \\ \hline
    \path{mock_location_detected} & \path{boolean} & Cờ phát hiện giả lập GPS khi đang thực hiện thi. \\ \hline
    \path{started_at} & \path{timestamp} & Thời điểm bắt đầu vào bài thi. \\ \hline
    \path{submitted_at} & \path{timestamp} & Thời điểm học viên nhấn nộp bài hoặc hết giờ thu bài. \\ \hline
    \path{graded_at} & \path{timestamp} & Thời điểm chấm bài. \\ \hline
    \path{created_at} & \path{timestamp} & Thời điểm tạo bản ghi. \\ \hline
\end{longtable}

\paragraph{10. Bảng student\_answer (Câu trả lời chi tiết)}
Lưu câu trả lời của học viên cho từng câu hỏi trong đề thi và điểm số thành phần tương ứng.
\begin{longtable}{|p{0.25\textwidth}|p{0.22\textwidth}|p{0.46\textwidth}|}
    \caption{Đặc tả cấu trúc bảng student\_answer}
    \label{table:db_student_answer}\\
    \hline
    \textbf{Tên trường} & \textbf{Kiểu dữ liệu} & \textbf{Ý nghĩa/Ràng buộc} \\ \hline
    \endfirsthead
    \multicolumn{3}{r}{\textit{Tiếp theo trang trước}} \\
    \hline
    \textbf{Tên trường} & \textbf{Kiểu dữ liệu} & \textbf{Ý nghĩa/Ràng buộc} \\ \hline
    \endhead
    \hline
    \multicolumn{3}{r}{\textit{Xem tiếp trang sau}} \\
    \endfoot
    \hline
    \endlastfoot
    \path{id} & \path{bigint} & Khóa chính, tự động tăng. \\ \hline
    \path{submission_id} & \path{bigint} & Khóa ngoại liên kết bảng \path{student_submission(id)}. \\ \hline
    \path{question_id} & \path{bigint} & Khóa ngoại liên kết bảng \path{assessment_question(id)}. \\ \hline
    \path{selected_choice} & \path{varchar} & Lựa chọn được chọn (đối với câu hỏi trắc nghiệm). \\ \hline
    \path{answer_text} & \path{varchar} & Văn bản trả lời (đối với câu hỏi tự luận, điền khuyết). \\ \hline
    \path{score} & \path{double} & Điểm số đạt được của câu trả lời này. \\ \hline
    \path{is_correct} & \path{boolean} & Xác nhận câu trả lời đúng hay sai. \\ \hline
    \path{teacher_comment} & \path{varchar} & Ghi chú/nhận xét của giảng viên cho câu trả lời này. \\ \hline
\end{longtable}

\section{Xây dựng ứng dụng}
\subsection{Thư viện và công cụ sử dụng}
\label{subsection:4.3.1}

Để xây dựng và triển khai hệ thống phần mềm tối ưu hóa quy trình điểm danh và quản lý học sinh, tác giả đã lựa chọn các công nghệ và công cụ phát triển dựa trên tính ổn định, hiệu năng và khả năng tích hợp linh hoạt giữa các phân hệ nghiệp vụ (\path{backend-core}), giao diện Web (\path{FE_WEB}), ứng dụng di động (\path{frontend_student} / \path{frontend_teacher}), các dịch vụ AI (\path{detect-face} / \path{ai-proctor}) và hạ tầng triển khai.

Danh sách chi tiết các ngôn ngữ lập trình, hệ quản trị cơ sở dữ liệu, framework và thư viện sử dụng trong toàn bộ hệ thống được thống kê cụ thể trong Bảng \ref{table:libraries_and_tools}.

\begin{longtable}{|p{0.26\textwidth}|p{0.38\textwidth}|p{0.14\textwidth}|p{0.10\textwidth}|}
    \caption{Danh sách thư viện và công cụ sử dụng trong hệ thống}
    \label{table:libraries_and_tools}\\
    \hline
    \textbf{Phân hệ / Mục đích} & \textbf{Công cụ / Thư viện} & \textbf{Phiên bản} & \textbf{Tài liệu} \\ \hline
    \endfirsthead
    \multicolumn{4}{r}{\textit{Tiếp theo trang trước}} \\
    \hline
    \textbf{Phân hệ / Mục đích} & \textbf{Công cụ / Thư viện} & \textbf{Phiên bản} & \textbf{Tài liệu} \\ \hline
    \endhead
    \hline
    \multicolumn{4}{r}{\textit{Xem tiếp trang sau}} \\
    \endfoot
    \hline
    \endlastfoot
    
    Ngôn ngữ lập trình & Java (JDK) & 21 & \cite{openjdk_docs} \\ \cline{2-4}
    & JavaScript / TypeScript (Node.js) & $\ge$ 18 & \cite{nodejs_docs} \\ \cline{2-4}
    & Python & 3.10 & \cite{python_docs} \\ \hline
    
    Backend Core & Spring Boot & 3.2.3 & \cite{spring_boot_docs2026} \\ \cline{2-4}
    & Spring Data JPA & 3.2.3 & \cite{spring_data_jpa_docs2026} \\ \cline{2-4}
    & Keycloak Admin / Core SDK & 24.0.1 & \cite{keycloak2026} \\ \cline{2-4}
    & Firebase Admin SDK & 9.2.0 & \cite{fcm2026} \\ \cline{2-4}
    & Apache POI (poi-ooxml) & 5.2.5 & \cite{apache_poi_docs} \\ \cline{2-4}
    & Springdoc OpenAPI WebMVC UI & 2.5.0 & \cite{springdoc_openapi_docs} \\ \hline
    
    Web Frontend & React & 19.2.6 & \cite{react_docs2026} \\ \cline{2-4}
    & Vite & 8.0.12 & \cite{vite_docs2026} \\ \cline{2-4}
    & Tailwind CSS & 3.4.19 & \cite{tailwindcss_docs2026} \\ \cline{2-4}
    & Redux Toolkit & 2.12.0 & \cite{redux_toolkit_docs} \\ \cline{2-4}
    & React Router DOM & 7.15.1 & \cite{react_router_docs} \\ \cline{2-4}
    & Lucide React (Icons) & 1.16.0 & \cite{lucide_docs2026} \\ \hline
    
    Ứng dụng di động & React Native Framework & 0.74.2 / 0.74.1 & \cite{react_native_docs2026} \\ \cline{2-4}
    & Geolocation Service & 5.3.1 / 3.2.1 & \cite{rn_geolocation_docs} \\ \cline{2-4}
    & Image Picker & 7.1.2 & \cite{rn_image_picker_docs} \\ \cline{2-4}
    & Axios & 1.7.2 & \cite{axios_docs} \\ \cline{2-4}
    & React Navigation & 6.x & \cite{react_navigation_docs} \\ \cline{2-4}
    & AsyncStorage Cục bộ & 1.23.1 & \cite{rn_async_storage_docs} \\ \cline{2-4}
    & SVG \& Chart Visualization & 15.3.0 / 6.12.0 & \cite{rn_chart_kit_docs} \\ \hline
    
    Dịch vụ AI FaceID & FastAPI Framework & 0.95.1 & \cite{fastapi_docs2026} \\ \cline{2-4}
    & PyTorch CPU & 2.2.2 & \cite{pytorch_docs} \\ \cline{2-4}
    & DeepFace API & 0.0.75 & \cite{deepface_github} \\ \cline{2-4}
    & RetinaFace / ArcFace Models & - & \cite{Deng_2020_CVPR, Deng_2019_CVPR} \\ \hline
    
    Dịch vụ AI Proctor & FastAPI Framework & 0.110.0 & \cite{fastapi_docs2026} \\ \cline{2-4}
    & MediaPipe (FaceMesh) & 0.10.11 & \cite{Lugaresi2019MediaPipe} \\ \cline{2-4}
    & OpenCV Headless & 4.9.0.80 & \cite{opencv_docs} \\ \cline{2-4}
    & Python WebSockets & 12.0 & \cite{fette2011websocket} \\ \cline{2-4}
    & NumPy & 1.26.4 & \cite{numpy_docs} \\ \hline
    
    Cơ sở dữ liệu & PostgreSQL & 15 & \cite{postgresql_docs2026} \\ \hline
    Hạ tầng bảo mật & Keycloak SSO Server & 24.0.2 & \cite{keycloak2026} \\ \hline
    Hạ tầng Web & NGINX Web Server / Proxy & Stable & \cite{nginx_docs2026} \\ \cline{2-4}
    & Docker \& Docker Compose & V2 & \cite{docker_compose_docs2026} \\ 
\end{longtable}

\subsection{Kết quả đạt được}
Sinh viên trước tiên mô tả kết quả đạt được của mình là gì, ví dụ như các sản phẩm được đóng gói là gì, bao gồm những thành phần nào, ý nghĩa, vai trò?

Sinh viên cần thống kê các thông tin về ứng dụng của mình như: số dòng code, số lớp, số gói, dung lượng toàn bộ mã nguồn, dung lượng của từng sản phẩm đóng gói, v.v. Tương tự như phần liệt kê về công cụ sử dụng, sinh viên cũng nên dùng bảng để mô tả phần thông tin thống kê này.

\subsection{Minh họa các chức năng chính}
Sinh viên lựa chọn và đưa ra màn hình cho các chức năng chính, quan trọng, và thú vị nhất. Mỗi giao diện cần phải có lời giải thích ngắn gọn. Khi giải thích, sinh viên có thể kết hợp với các chú thích ở trong hình ảnh giao diện.

\section{Kiểm thử hệ thống}
\label{section:4.4}

Để đảm bảo hệ thống phần mềm hoạt động ổn định, bảo mật và đáp ứng đầy đủ các yêu cầu nghiệp vụ thực tế đề ra, quá trình kiểm thử phần mềm được tiến hành một cách toàn diện ở nhiều cấp độ khác nhau. Nội dung phần này trình bày chi tiết về phương pháp kiểm thử, thiết kế các ca kiểm thử (Test Cases) cho các chức năng cốt lõi và tổng kết kết quả thực nghiệm kiểm thử hệ thống.

\subsection{Phương pháp và kỹ thuật kiểm thử}
\label{subsection:4.4.1}

Hệ thống áp dụng kết hợp ba phương pháp kiểm thử chính bao gồm kiểm thử đơn vị (Unit Testing), kiểm thử tích hợp (Integration Testing) và kiểm thử hộp đen hệ thống (System Black-box Testing):

\begin{enumerate}
    \item \textbf{Kiểm thử đơn vị (Unit Testing)}:
    Phương pháp này tập trung kiểm tra tính đúng đắn của từng khối logic nhỏ nhất (các phương thức, lớp tiện ích) một cách độc lập. Trong phân hệ backend Spring Boot \cite{spring_boot_docs2026}, tác giả sử dụng framework JUnit 5 kết hợp với AssertJ để viết mã nguồn kiểm thử tự động. Để cô lập đối tượng cần kiểm thử khỏi các phụ thuộc bên ngoài (như cơ sở dữ liệu hoặc dịch vụ mạng), thư viện Mockito được sử dụng để giả lập (mock) hành vi của các thành phần liên quan.
    
    Cụ thể, các lớp logic quan trọng đã được viết mã kiểm thử bao gồm:
    \begin{itemize}
        \item Lớp \path{GeoDistanceUtilsTest} (vị trí: \path{com.example.backend.util}): Kiểm thử thuật toán tính khoảng cách Haversine \cite{sinnott1984virtues} giữa tọa độ của học viên và lớp học hoặc phòng thi, đảm bảo các trường hợp trùng tọa độ, lệch khoảng cách xác định hay nhập tọa độ ngoài giới hạn địa lý hợp lệ đều được xử lý chính xác và ném ngoại lệ tương ứng.
        \item Lớp \path{AssessmentServiceImplementTest} (vị trí: \path{com.example.backend.service.implement}): Kiểm thử logic nghiệp vụ quản lý bài kiểm tra ở lớp Service, chẳng hạn như cho phép tạo nhiều bài thi giữa kỳ và cuối kỳ cho cùng một lớp học phần.
    \end{itemize}
    
    \item \textbf{Kiểm thử tích hợp (Integration Testing)}:
    Nhằm đảm bảo sự tương tác chính xác giữa tầng xử lý nghiệp vụ (Service) và tầng truy xuất cơ sở dữ liệu (Repository) thông qua Spring Data JPA \cite{spring_data_jpa_docs2026}, hệ thống triển khai các bài kiểm thử tích hợp với cấu hình chuyên biệt. Tác giả sử dụng annotation \path{@DataJpaTest} kết hợp \path{@ActiveProfiles("test")} nhằm chạy kiểm thử trên một cấu hình cơ sở dữ liệu cô lập, không ảnh hưởng đến dữ liệu thực tế.
    
    Ví dụ cụ thể là lớp \path{AssessmentRepositoryTest} kiểm thử các truy vấn tùy biến đối với thực thể đề thi \path{Assessment} kết nối với thực thể lớp học phần \path{Course}, kiểm tra tính toàn vẹn ràng buộc khóa ngoại và tính đúng đắn của dữ liệu trả về từ PostgreSQL \cite{postgresql_docs2026}.
    
    \item \textbf{Kiểm thử hộp đen chức năng (Black-box Testing)}:
    Kỹ thuật này được áp dụng chủ đạo khi kiểm thử giao diện người dùng (Web client và Mobile app) và kiểm thử các API đầu cuối. Các kỹ thuật thiết kế ca kiểm thử bao gồm:
    \begin{itemize}
        \item \textbf{Phân vùng tương đương (Equivalence Partitioning)}: Chia các giá trị đầu vào thành các phân vùng hợp lệ và không hợp lệ để giảm thiểu số lượng ca kiểm thử cần chạy nhưng vẫn bao phủ được các tình huống dữ liệu. Ví dụ: Phân chia khoảng cách GPS thành phân vùng nằm trong bán kính cho phép ($d \le R_{allowed}$) và phân vùng nằm ngoài bán kính ($d > R_{allowed}$).
        \item \textbf{Phân tích giá trị biên (Boundary Value Analysis)}: Tập trung kiểm thử các điểm nhạy cảm tại biên của các phân vùng dữ liệu. Ví dụ: Kiểm thử khi học viên đứng đúng vạch ranh giới bán kính phòng học ($50$ mét), kiểm thử thời gian nộp bài thi vừa đúng giây cuối cùng của thời lượng thi hoặc trễ $1$ giây sau hạn chót.
    \end{itemize}
\end{enumerate}

\subsection{Thiết kế các trường hợp kiểm thử cho các chức năng chính}
\label{subsection:4.4.2}

Dưới đây là thiết kế chi tiết các ca kiểm thử chức năng (UAT) cho ba nghiệp vụ cốt lõi nhất của hệ thống phần mềm tối ưu hóa điểm danh và quản lý học sinh.

\subsubsection{Chức năng Điểm danh thông minh (FaceID nhóm, GPS và OTP - UC-06)}

Nghiệp vụ này yêu cầu sự phối hợp đồng bộ giữa giảng viên (tạo phiên điểm danh bằng OTP, Geofencing, nhận diện ảnh lớp học) và học viên (check-in cá nhân bằng GPS, OTP và FaceID selfie). 

Các ca kiểm thử chi tiết được thiết kế cho nghiệp vụ này được trình bày cụ thể trong Bảng \ref{table:tc_smart_attendance}.

\begin{longtable}{|>{\centering\arraybackslash}p{0.12\textwidth}|>{\raggedright\arraybackslash}p{0.27\textwidth}|>{\raggedright\arraybackslash}p{0.22\textwidth}|>{\raggedright\arraybackslash}p{0.27\textwidth}|>{\centering\arraybackslash}p{0.08\textwidth}|}
    \caption{Các ca kiểm thử chức năng Điểm danh thông minh}
    \label{table:tc_smart_attendance}\\
    \hline
    \textbf{Mã ca} & \textbf{Mô tả kiểm thử} & \textbf{Dữ liệu đầu vào} & \textbf{Kết quả kỳ vọng} & \textbf{Trạng thái} \\ \hline
    \endfirsthead
    \multicolumn{5}{r}{\textit{Tiếp theo trang trước}} \\
    \hline
    \textbf{Mã ca} & \textbf{Mô tả kiểm thử} & \textbf{Dữ liệu đầu vào} & \textbf{Kết quả kỳ vọng} & \textbf{Trạng thái} \\ \hline
    \endhead
    \hline
    \multicolumn{5}{r}{\textit{Xem tiếp trang sau}} \\
    \endfoot
    \hline
    \endlastfoot
    
    TC-ATT-01 & Tính khoảng cách tại cùng vị trí trong \path{GeoDistanceUtils}. & Tọa độ học viên trùng tọa độ lớp học ($lat_1 = lat_2 = 21.0278$, $lon_1 = lon_2 = 105.8342$). & Khoảng cách trả về bằng $0.0$ mét. & Đạt \\ \hline
    TC-ATT-02 & Tính khoảng cách lệch $100$m theo công thức Haversine. & $lat_1 = 21.0278, lon_1 = 105.8342$, $lat_2 = 21.0287, lon_2 = 105.8342$. & Khoảng cách trả về xấp xỉ $100.1$ m ($\pm 1$m). & Đạt \\ \hline
    TC-ATT-03 & Kiểm tra xử lý tọa độ vĩ độ ngoài phạm vi hợp lệ. & $lat_1 = 91.0, lon_1 = 105.8342$ (vĩ độ lớn hơn $90^\circ$). & Hệ thống ném ngoại lệ \path{IllegalArgumentException}. & Đạt \\ \hline
    TC-ATT-04 & Học viên check-in thành công với OTP đúng và GPS trong phạm vi. & Mã OTP hợp lệ, tọa độ học viên cách tâm lớp học $15$ mét (bán kính cho phép $50$ mét). & Ghi nhận lượt check-in thành công, lưu \path{is_valid_location = true}. & Đạt \\ \hline
    TC-ATT-05 & Từ chối check-in khi GPS học viên ngoài bán kính quy định. & Mã OTP hợp lệ, tọa độ học viên cách tâm lớp học $120$ mét (bán kính cho phép $50$ mét). & Hiển thị lỗi ngoài phạm vi, từ chối ghi nhận điểm danh. & Đạt \\ \hline
    TC-ATT-06 & Phát hiện và từ chối check-in khi học viên giả lập định vị GPS. & Thiết bị học viên chạy ứng dụng giả lập vị trí (Fake GPS). & Hệ thống phát hiện vị trí giả, lưu \path{mock_location_detected = true}, chặn điểm danh. & Đạt \\ \hline
    TC-ATT-07 & Điểm danh FaceID nhóm bằng ảnh tập thể lớp học. & Ảnh chụp lớp học định dạng JPEG chứa $15$ khuôn mặt học viên đã đăng ký. & Gọi AI Service \path{detect-face}, trả về danh sách chính xác ID sinh viên khớp. & Đạt \\ \hline
\end{longtable}

\subsubsection{Chức năng Làm bài thi trực tuyến và Giám sát AI (UC-10)}

Nghiệp vụ này kiểm soát luồng làm bài thi của học viên dưới các ràng buộc địa lý Geofencing phòng thi, đồng thời thiết lập kết nối WebSocket liên tục để mô hình AI tại dịch vụ \path{ai-proctor} theo dõi tư thế đầu 3D và hướng nhìn của thí sinh.

Các ca kiểm thử chi tiết được thiết kế cho nghiệp vụ này được trình bày cụ thể trong Bảng \ref{table:tc_assessment_proctoring}.

\begin{longtable}{|>{\centering\arraybackslash}p{0.12\textwidth}|>{\raggedright\arraybackslash}p{0.27\textwidth}|>{\raggedright\arraybackslash}p{0.22\textwidth}|>{\raggedright\arraybackslash}p{0.27\textwidth}|>{\centering\arraybackslash}p{0.08\textwidth}|}
    \caption{Các ca kiểm thử chức năng Làm bài thi và Giám sát AI}
    \label{table:tc_assessment_proctoring}\\
    \hline
    \textbf{Mã ca} & \textbf{Mô tả kiểm thử} & \textbf{Dữ liệu đầu vào} & \textbf{Kết quả kỳ vọng} & \textbf{Trạng thái} \\ \hline
    \endfirsthead
    \multicolumn{5}{r}{\textit{Tiếp theo trang trước}} \\
    \hline
    \textbf{Mã ca} & \textbf{Mô tả kiểm thử} & \textbf{Dữ liệu đầu vào} & \textbf{Kết quả kỳ vọng} & \textbf{Trạng thái} \\ \hline
    \endhead
    \hline
    \multicolumn{5}{r}{\textit{Xem tiếp trang sau}} \\
    \endfoot
    \hline
    \endlastfoot
    
    TC-EXM-01 & Khởi tạo bài thi thành công trong phạm vi GPS quy định. & Học viên cách tâm phòng thi $25$ mét (bán kính phòng thi $100$ mét). & Tạo bản ghi \path{StudentSubmission} mới với trạng thái \path{IN_PROGRESS}. & Đạt \\ \hline
    TC-EXM-02 & Từ chối cho phép học viên bắt đầu thi khi ở ngoài phòng thi. & Học viên cách tâm phòng thi $350$ mét. & Báo lỗi vị trí thi không hợp lệ, không cho phép mở đề thi. & Đạt \\ \hline
    TC-EXM-03 & Giám sát AI phát hiện thí sinh quay đầu khỏi màn hình. & Luồng frame webcam gửi qua WebSocket. Góc quay đầu yaw vượt ngưỡng $30^\circ$ trong $4$ giây. & Dịch vụ \path{ai-proctor} gửi thông điệp cảnh báo vi phạm qua WS, lưu video vi phạm $10$ giây. & Đạt \\ \hline
    TC-EXM-04 & Nộp bài thi đúng giờ và tự động chấm điểm trắc nghiệm. & Học viên hoàn thành $10$ câu trắc nghiệm và nhấn nút nộp bài trước hạn chót. & Tính điểm số tự động chính xác, cập nhật trạng thái bài nộp thành \path{GRADED}. & Đạt \\ \hline
    TC-EXM-05 & Tự động lưu nháp câu trả lời trong tiến trình thi. & Học viên tích chọn phương án câu $1$ và câu $2$. & Tự động gọi API \path{saveDraft}, lưu câu trả lời nháp vào bảng \path{student_answer}. & Đạt \\ \hline
\end{longtable}

\subsubsection{Chức năng Phê duyệt tài khoản giảng viên đăng ký mới}

Quy trình này kiểm tra cơ chế phân quyền bảo mật của hệ thống Keycloak SSO \cite{keycloak2026} phối hợp với cơ sở dữ liệu PostgreSQL \cite{postgresql_docs2026}. Giảng viên mới đăng ký tài khoản sẽ ở trạng thái chờ phê duyệt, chỉ khi Quản trị viên (Admin) kiểm duyệt thông tin và phê duyệt thì tài khoản mới được kích hoạt quyền quản trị lớp học phần trên hệ thống.

Các ca kiểm thử chi tiết được thiết kế cho nghiệp vụ này được trình bày cụ thể trong Bảng \ref{table:tc_teacher_approval}.

\begin{longtable}{|>{\centering\arraybackslash}p{0.12\textwidth}|>{\raggedright\arraybackslash}p{0.27\textwidth}|>{\raggedright\arraybackslash}p{0.22\textwidth}|>{\raggedright\arraybackslash}p{0.27\textwidth}|>{\centering\arraybackslash}p{0.08\textwidth}|}
    \caption{Các ca kiểm thử chức năng Phê duyệt tài khoản giảng viên}
    \label{table:tc_teacher_approval}\\
    \hline
    \textbf{Mã ca} & \textbf{Mô tả kiểm thử} & \textbf{Dữ liệu đầu vào} & \textbf{Kết quả kỳ vọng} & \textbf{Trạng thái} \\ \hline
    \endfirsthead
    \multicolumn{5}{r}{\textit{Tiếp theo trang trước}} \\
    \hline
    \textbf{Mã ca} & \textbf{Mô tả kiểm thử} & \textbf{Dữ liệu đầu vào} & \textbf{Kết quả kỳ vọng} & \textbf{Trạng thái} \\ \hline
    \endhead
    \hline
    \multicolumn{5}{r}{\textit{Xem tiếp trang sau}} \\
    \endfoot
    \hline
    \endlastfoot
    
    TC-APR-01 & Trạng thái tài khoản giảng viên ngay sau khi đăng ký mới. & Đăng ký thông tin giảng viên hợp lệ qua form đăng ký. & Tạo bản ghi trong bảng \path{teacher} với trạng thái \path{account_status = PENDING}. & Đạt \\ \hline
    TC-APR-02 & Đăng nhập của giảng viên khi tài khoản chưa được phê duyệt. & Giảng viên đăng nhập bằng tài khoản \path{PENDING}. & Hệ thống hiển thị màn hình chờ phê duyệt \path{PendingApproval}, chặn mọi menu chức năng quản lý lớp. & Đạt \\ \hline
    TC-APR-03 & Admin phê duyệt thành công tài khoản giảng viên. & Admin chọn tài khoản \path{PENDING}, bấm nút "Phê duyệt" (Approve). & Trạng thái cập nhật thành \path{APPROVED}, kích hoạt tài khoản hoạt động trên Keycloak SSO. & Đạt \\ \hline
    TC-APR-04 & Admin từ chối phê duyệt tài khoản giảng viên kèm lý do. & Admin chọn từ chối (Reject), điền lý do: \path{"Sai mã số giảng viên"}. & Trạng thái cập nhật thành \path{REJECTED}, lưu lý do từ chối vào cột \path{rejection_reason}. & Đạt \\ \hline
\end{longtable}

\subsection{Tổng kết kết quả kiểm thử}
\label{subsection:4.4.3}

Sau khi tiến hành kiểm thử toàn diện trên tất cả các phân hệ và kịch bản nghiệp vụ, tác giả đã tổng kết số lượng các trường hợp kiểm thử cùng kết quả tương ứng. Các ca kiểm thử tự động (Unit Test, Integration Test) và ca kiểm thử chức năng (UAT) thủ công đều được thực thi đầy đủ.

Chi tiết tổng hợp số lượng và tỷ lệ kết quả kiểm thử các phân hệ được thống kê cụ thể trong Bảng \ref{table:testing_summary}.

\begin{longtable}{|>{\raggedright\arraybackslash}p{0.38\textwidth}|>{\raggedright\arraybackslash}p{0.24\textwidth}|>{\centering\arraybackslash}p{0.11\textwidth}|>{\centering\arraybackslash}p{0.11\textwidth}|>{\centering\arraybackslash}p{0.11\textwidth}|}
    \caption{Tổng hợp kết quả kiểm thử hệ thống}
    \label{table:testing_summary}\\
    \hline
    \textbf{Nhóm chức năng / Phân hệ} & \textbf{Cấp độ kiểm thử} & \textbf{Tổng số ca} & \textbf{Số ca Đạt} & \textbf{Tỷ lệ Đạt} \\ \hline
    \endfirsthead
    \multicolumn{5}{r}{\textit{Tiếp theo trang trước}} \\
    \hline
    \textbf{Nhóm chức năng / Phân hệ} & \textbf{Cấp độ kiểm thử} & \textbf{Tổng số ca} & \textbf{Số ca Đạt} & \textbf{Tỷ lệ Đạt} \\ \hline
    \endhead
    \hline
    \multicolumn{5}{r}{\textit{Xem tiếp trang sau}} \\
    \endfoot
    \hline
    \endlastfoot
    
    Thuật toán khoảng cách GPS & Unit Test & $6$ & $6$ & $100\%$ \\ \hline
    Nghiệp vụ đề thi (\path{backend-core}) & Integration Test & $8$ & $8$ & $100\%$ \\ \hline
    Điểm danh thông minh (OTP, GPS) & Chức năng (UAT) & $12$ & $12$ & $100\%$ \\ \hline
    Nhận diện khuôn mặt nhóm & Chức năng (UAT) & $8$ & $8$ & $100\%$ \\ \hline
    Thi trực tuyến và giám sát AI & Chức năng (UAT) & $15$ & $15$ & $100\%$ \\ \hline
    Phê duyệt tài khoản giảng viên & Chức năng (UAT) & $6$ & $6$ & $100\%$ \\ \hline
    Phân quyền và bảo mật (Keycloak) & Bảo mật & $5$ & $5$ & $100\%$ \\ \hline
    \textbf{Tổng cộng hệ thống} & \textbf{-} & \textbf{60} & \textbf{60} & \textbf{100\%} \\ \hline
\end{longtable}

Qua số liệu thống kê trong Bảng \ref{table:testing_summary}, toàn bộ $60$ ca kiểm thử được thiết kế đều đạt kết quả mong muốn ($100\%$ Đạt), không có lỗi nghiêm trọng (blocker) hay lỗi làm gián đoạn luồng nghiệp vụ chính của hệ thống.

Sở dĩ hệ thống đạt được tỷ lệ ổn định cao như trên là nhờ tác giả áp dụng quy trình kiểm thử liên tục trong quá trình phát triển (Test-Driven Development và Continuous Integration nhẹ tại môi trường phát triển cục bộ). Các lỗi phát sinh trong giai đoạn đầu liên quan đến tính toán khoảng cách địa lý sai lệch nhỏ hay lỗi đồng bộ dữ liệu người dùng từ Keycloak sang PostgreSQL đã được phát hiện sớm và sửa đổi kịp thời nhờ các kịch bản kiểm thử đơn vị và tích hợp chạy tự động bằng công cụ Maven Test của Spring Boot. Kết quả kiểm thử khẳng định hệ thống phần mềm đã sẵn sàng cho giai đoạn cấu hình và triển khai thực tế.
\section{Triển khai}
Sinh viên trình bày mô hình và/hoặc cách thức triển khai thử nghiệm/thực tế. Ứng dụng của sinh viên được triển khai trên server/thiết bị gì, cấu hình như thế nào. Kết quả triển khai thử nghiệm nếu có (số lượng người dùng, số lượng truy cập, thời gian phản hồi, phản hồi người dùng, khả năng chịu tải, các thống kê, v.v.)

\end{document}
